\documentclass[answers]{exam}
\usepackage{../../../../mypackages}
\usepackage{../../../../macros}

\tcbuselibrary{theorems,skins,hooks}

\definecolor{myindbg}{HTML}{DFF5EC}
\definecolor{myindfr}{HTML}{1F6F4A}

\newtcolorbox{Indication}{
  enhanced,
  colback=myindbg,
  colframe=myindfr,
  arc=6mm,
  rounded corners,
  boxrule=0.6mm,
  left=3mm,
  right=3mm,
  top=2mm,
  bottom=2mm
}

\title{Correction exercices chapitre sur angles}
\author{N. Bancel}
\date{15 février 2026}

% Mise en forme des solutions en bleu
\SolutionEmphasis{\color{blue}}
\renewcommand{\solutiontitle}{\noindent}

\begin{document}

\dsheader{Collège Lycée Suger}{Mathématiques}{Année 2025-2026}{6ème}
{\let\newpage\relax\maketitle}

\section*{Corrigé}

%------------------------------------------------
\section{Exercice 99 : Calculer l'angle $\widehat{EGF}$}
\fig{0.6}{./angle_exo_2.jpg}{Exercice 99}

\begin{solution}
On cherche la mesure de l'angle $\widehat{EGF}$ (angle au point $G$).

\begin{Indication}
Les arcs verts indiquent des angles égaux : \\
\begin{compactitem}
\item au point $D$, $[DB]$ est la bissectrice de l'angle $\widehat{CDE}$ ;
\item au point $F$, $[FE]$ est la bissectrice de l'angle $\widehat{DFG}$.
\end{compactitem}
\end{Indication}

\begin{compactenum}
\item \textbf{Calcul de l'angle $\widehat{ADB}$.}

On sait que les points $A$, $C$ et $D$ sont alignés, donc $[AD]$ est le même alignement que $[AC]$.

À $B$, on lit :
\[
\widehat{ABD} = 37^\circ + 31^\circ = 68^\circ
\]
À $A$, on lit :
\[
\widehat{BAD} = 68^\circ
\]

On sait que dans le triangle $ABD$ :
\begin{align*}
\widehat{ABD} + \widehat{BAD} + \widehat{ADB} &= 180^\circ \\
\widehat{ADB} &= 180^\circ - \widehat{ABD} - \widehat{BAD}
\end{align*}
Application numérique :
\begin{align*}
\widehat{ADB} &= 180^\circ - 68^\circ - 68^\circ \\
\widehat{ADB} &= 44^\circ
\end{align*}

\item \textbf{Utilisation de la bissectrice au point $D$.}

On sait que $[DB]$ est la bissectrice de l'angle $\widehat{CDE}$.

Or, comme $A$, $C$ et $D$ sont alignés, l'angle $\widehat{CDB}$ est le même que l'angle $\widehat{ADB}$ :
\[
\widehat{CDB} = \widehat{ADB} = 44^\circ
\]

Donc, par définition de la bissectrice :
\[
\widehat{BDE} = \widehat{CDB} = 44^\circ
\]

\item \textbf{Calcul de l'angle $\widehat{BED}$.}

On sait que les points $B$, $E$ et $F$ sont alignés, donc $[BE]$ est une droite passant par $E$.

À $B$, on lit :
\[
\widehat{DBE} = 32^\circ
\]

Dans le triangle $BDE$ :
\begin{align*}
\widehat{DBE} + \widehat{BDE} + \widehat{BED} &= 180^\circ \\
\widehat{BED} &= 180^\circ - \widehat{DBE} - \widehat{BDE}
\end{align*}
Application numérique :
\begin{align*}
\widehat{BED} &= 180^\circ - 32^\circ - 44^\circ \\
\widehat{BED} &= 104^\circ
\end{align*}

\item \textbf{Passage à l'angle $\widehat{DEF}$.}

On sait que $B$, $E$ et $F$ sont alignés, donc les demi-droites $[EB)$ et $[EF)$ sont opposées.

Or, des angles supplémentaires sur une même droite vérifient :
\[
\widehat{BED} + \widehat{DEF} = 180^\circ
\]
Donc :
\begin{align*}
\widehat{DEF} &= 180^\circ - \widehat{BED} \\
\widehat{DEF} &= 180^\circ - 104^\circ \\
\widehat{DEF} &= 76^\circ
\end{align*}

\item \textbf{Calcul de l'angle $\widehat{DFE}$ dans le triangle $DEF$.}

On lit sur la figure :
\[
\widehat{FDE} = 60^\circ
\]
et on vient de trouver :
\[
\widehat{DEF} = 76^\circ
\]

Dans le triangle $DEF$ :
\begin{align*}
\widehat{FDE} + \widehat{DEF} + \widehat{DFE} &= 180^\circ \\
\widehat{DFE} &= 180^\circ - \widehat{FDE} - \widehat{DEF}
\end{align*}
Application numérique :
\begin{align*}
\widehat{DFE} &= 180^\circ - 60^\circ - 76^\circ \\
\widehat{DFE} &= 44^\circ
\end{align*}

\item \textbf{Utilisation de la bissectrice au point $F$ puis calcul de $\widehat{EGF}$.}

On sait que $[FE]$ est la bissectrice de l'angle $\widehat{DFG}$.

Or, cela signifie :
\[
\widehat{DFE} = \widehat{EFG}
\]
Donc :
\[
\widehat{EFG} = 44^\circ
\]

Dans le triangle $EFG$, on lit :
\[
\widehat{FEG} = 70^\circ
\]
Alors :
\begin{align*}
\widehat{EFG} + \widehat{FEG} + \widehat{EGF} &= 180^\circ \\
\widehat{EGF} &= 180^\circ - \widehat{EFG} - \widehat{FEG}
\end{align*}
Application numérique :
\begin{align*}
\widehat{EGF} &= 180^\circ - 44^\circ - 70^\circ \\
\widehat{EGF} &= 66^\circ
\end{align*}

\end{compactenum}

\[
\boxed{\widehat{EGF} = 66^\circ}
\]
\end{solution}

%------------------------------------------------
\section{Exercice 94 : Angles et triangle}
\fig{0.6}{./angles_exo_1.jpg}{Exercice 94}

\begin{solution}
Sur la figure, les points $A$, $E$ et $M$ sont alignés.

Les marques d'égalité indiquent :
\[
AF = AE \quad \text{et} \quad AE = EF
\]
Donc :
\[
AF = AE = EF
\]
Le carré au point $F$ indique :
\[
(AF) \perp (FM)
\]

\begin{parts}

%------------------------
\part Quelle est la mesure de l'angle $\widehat{AEF}$ ? Justifier.

On sait que $AF = AE = EF$.

Or, dans un triangle qui a ses \textbf{trois côtés égaux}, les trois angles sont égaux, donc chacun vaut $60^\circ$.

Donc :
\[
\boxed{\widehat{AEF} = 60^\circ}
\]

%------------------------
\part En déduire la mesure de l'angle $\widehat{FEM}$.

On sait que $A$, $E$ et $M$ sont alignés.

Or, si $A$, $E$, $M$ sont alignés, alors les demi-droites $[EA)$ et $[EM)$ sont opposées, donc les angles $\widehat{AEF}$ et $\widehat{FEM}$ sont supplémentaires :
\[
\widehat{AEF} + \widehat{FEM} = 180^\circ
\]
Donc :
\begin{align*}
\widehat{FEM} &= 180^\circ - \widehat{AEF} \\
\widehat{FEM} &= 180^\circ - 60^\circ \\
\widehat{FEM} &= 120^\circ
\end{align*}

\[
\boxed{\widehat{FEM} = 120^\circ}
\]

%------------------------
\part Calculer la mesure de l'angle $\widehat{EMF}$.

On commence par calculer $\widehat{EFM}$.

On sait que $(AF) \perp (FM)$, donc :
\[
\widehat{AFM} = 90^\circ
\]
On sait aussi que dans le triangle $AEF$ (équilatéral), l'angle au point $F$ vaut :
\[
\widehat{AFE} = 60^\circ
\]
Or, la demi-droite $[FE)$ est à l'intérieur de l'angle $\widehat{AFM}$, donc :
\[
\widehat{AFM} = \widehat{AFE} + \widehat{EFM}
\]
Donc :
\begin{align*}
\widehat{EFM} &= \widehat{AFM} - \widehat{AFE} \\
\widehat{EFM} &= 90^\circ - 60^\circ \\
\widehat{EFM} &= 30^\circ
\end{align*}

Dans le triangle $EFM$ :
\begin{align*}
\widehat{FEM} + \widehat{EFM} + \widehat{EMF} &= 180^\circ \\
\widehat{EMF} &= 180^\circ - \widehat{FEM} - \widehat{EFM}
\end{align*}
Application numérique :
\begin{align*}
\widehat{EMF} &= 180^\circ - 120^\circ - 30^\circ \\
\widehat{EMF} &= 30^\circ
\end{align*}

\[
\boxed{\widehat{EMF} = 30^\circ}
\]

%------------------------
\part Reproduire la figure dans le cas où $AF = 4$ cm.

\begin{compactenum}
\item Tracer le segment $[AF]$ de longueur $\SI{4}{cm}$.
\item En $F$, tracer la droite $(FM)$ perpendiculaire à $(AF)$.
\item Tracer le cercle de centre $A$ et de rayon $\SI{4}{cm}$.
\item Tracer le cercle de centre $F$ et de rayon $\SI{4}{cm}$.
\item Leur intersection donne le point $E$ (on choisit le point situé du côté du segment $[FM]$).
\item Tracer la droite $(AE)$ ; elle coupe la droite $(FM)$ en $M$ (ainsi $A$, $E$, $M$ sont alignés).
\end{compactenum}

%------------------------
\part Construire la perpendiculaire à $[FM]$ passant par $E$. Elle coupe $[FM]$ en $P$.

Tracer par $E$ la droite perpendiculaire à $(FM)$ : elle coupe le segment $[FM]$ au point $P$.

%------------------------
\part Calculer la mesure de l'angle $\widehat{PEM}$.

On sait que $(EP) \perp (FM)$, donc l'angle entre $(EP)$ et $(FM)$ vaut $90^\circ$.

Or, l'angle entre $(EM)$ et $(FM)$ est l'angle $\widehat{EMF}$ (puisque $M$ est sur $(FM)$), et on a trouvé :
\[
\widehat{EMF} = 30^\circ
\]

Donc l'angle entre $(EP)$ et $(EM)$ vaut :
\begin{align*}
\widehat{PEM} &= 90^\circ - \widehat{EMF} \\
\widehat{PEM} &= 90^\circ - 30^\circ \\
\widehat{PEM} &= 60^\circ
\end{align*}

\[
\boxed{\widehat{PEM} = 60^\circ}
\]

\end{parts}
\end{solution}

\end{document}
